\documentclass[11pt,a4paper]{report}

\usepackage[utf8]{inputenc}
\usepackage[margin=1in]{geometry}
\usepackage{booktabs}
\usepackage{hyperref}
\usepackage{amsmath,amssymb}
\usepackage{graphicx}
\usepackage{listings}
\usepackage{xcolor}
\usepackage{titlesec}
\usepackage{fancyhdr}

\hypersetup{
    colorlinks=true,
    linkcolor=blue,
    filecolor=magenta,      
    urlcolor=cyan,
    citecolor=blue,
}

\lstset{
    basicstyle=\ttfamily\small,
    backgroundcolor=\color{gray!10},
    breaklines=true,
    frame=single,
    numbers=left,
    numberstyle=\tiny\color{gray},
    keywordstyle=\color{blue}\bfseries,
    stringstyle=\color{teal},
    commentstyle=\color{gray}\itshape,
    showstringspaces=false
}

\setlength{\headheight}{14pt}
\pagestyle{fancy}
\fancyhf{}
\rhead{\textbf{HELIOS Reference Manual}}
\lhead{Heavy Ion Linac Intelligent Orchestration}
\rfoot{Page \thepage}

\title{
    \vspace{2cm}
    \Huge \textbf{HELIOS Framework} \\
    \Large \textbf{User \& Developer Reference Manual} \\
    \vspace{0.5cm}
    \large Heavy Ion Linac Intelligent Orchestration \\
    \normalsize \textit{Autonomous, Safety-First Accelerator Control \& Commissioning Platform}
    \vspace{2cm}
}
\author{\textbf{HELIOS Collaboration} \\ Chong Shik Park (\texttt{cspark7701@gmail.com})}
\date{\today}

\begin{document}

\maketitle
\tableofcontents
\newpage

\chapter{System Overview and Architecture}

\section{Introduction}
\textbf{HELIOS} (\textbf{HE}avy ion \textbf{L}inac \textbf{I}ntelligent \textbf{O}rchestration) is a distributed multi-agent artificial intelligence framework engineered for commissioning, operating, and optimizing heavy-ion linear accelerators (LEBT, RFQ, MEBT, SCL, and HEBT).

Aligned directly with the \textbf{Plan-First} paradigm established by the Lawrence Berkeley National Laboratory (LBNL) Osprey architecture, HELIOS enforces strict separation of concerns between optimization heuristics, digital twin simulation validations, deterministic safety gatekeepers, and hardware control layers.

\section{Core Architectural Principles}
\begin{enumerate}
    \item \textbf{Pre-Flight Simulation Verification}: No setpoint is ever delivered to live hardware without prior validation through high-fidelity codes (ANL TRACK via Wine, LBNL IMPACT-Z via \texttt{lume-impact}) or sub-millisecond neural surrogates.
    \item \textbf{Deterministic Safety-First Actuation}: Optimization algorithms do not talk directly to EPICS. All actuations pass through rate-of-change limiters, slew rate clamps, and the \texttt{SafetyAgent} rollback stack.
    \item \textbf{Plan-First Execution Lifecycle}: High-level natural language goals are first compiled into structured, reproducible execution plans:
    \[
    \text{DRAFT} \longrightarrow \text{SIMULATION\_VERIFIED} \longrightarrow \text{SAFETY\_APPROVED} \longrightarrow \text{EXECUTING} \longrightarrow \text{COMPLETED / ROLLED\_BACK}
    \]
    \item \textbf{Dynamic Capability Classification (Osprey Pattern)}: Instead of injecting hundreds of tools into the LLM context, HELIOS dynamically filters capabilities based on the operator's intent and domain keywords, keeping prompt sizes compact and execution deterministic.
    \item \textbf{Graduated Autonomy}: Enforces clear human-in-the-loop sign-offs graduated from Advisory (Level 0) to Fully Autonomous (Level 4).
\end{enumerate}

\section{Agent Taxonomy and Roles}
The framework partitions operational responsibilities among specialized autonomous agents:
\begin{enumerate}
    \item \textbf{Master Orchestrator Agent (\texttt{HELIOS-Orchestrator})}: Decomposes high-level physicist directives into formal execution plans, manages multi-agent coordination, and interfaces with human operators.
    \item \textbf{Physics and Simulation Agent (\texttt{HELIOS-Physics})}: Manages the digital twin model, executes ANL TRACK and LBNL IMPACT-Z simulations, and evaluates deep neural surrogate models.
    \item \textbf{Diagnostics and Telemetry Agent (\texttt{HELIOS-Diagnostics})}: Ingests real-time signals from Beam Position Monitors (BPMs), Beam Loss Monitors (BLMs), and wire scanners; evaluates beam transmission and detects operational anomalies.
    \item \textbf{Tuning and Optimization Agent (\texttt{HELIOS-Tuner})}: Executes algorithmic optimization including response-matrix SVD orbit correction, multi-objective Bayesian Optimization, and Reinforcement Learning continuous control.
    \item \textbf{Safety and Verification Agent (\texttt{HELIOS-Safety})}: Serves as a deterministic gatekeeper enforcing hard software limits, slew-rate constraints, cryo-quench margins, and instant rollback stacks.
    \item \textbf{Control and Hardware Interface Agent (\texttt{HELIOS-Control})}: Interfaces with the live control system via EPICS Channel Access (CA) and pvAccess (PVA) with atomic readback verification.
\end{enumerate}

\section{Graduated Autonomy Levels}
HELIOS implements five graduated operational autonomy levels for human-in-the-loop gating:
\begin{table}[h]
\centering
\begin{tabular}{ccll}
\toprule
\textbf{Level} & \textbf{Name} & \textbf{Description} & \textbf{Operator Role} \\
\midrule
0 & Advisory & Passive telemetry monitoring \& anomaly alerts & Manual action only \\
1 & Assisted & Pre-flight verified plans generated & Mandatory approval per setpoint \\
2 & Supervised & Automated tuning within tight envelopes & Supervisory abort authority \\
3 & Bounded Auto & Closed-loop optimization within boundaries & Passive monitoring \\
4 & Fully Auto & Self-driving linac domain commissioning & Goal-setting oversight only \\
\bottomrule
\end{tabular}
\caption{Operational Autonomy Levels in HELIOS.}
\end{table}

\chapter{Installation and Environment Setup}

\section{Prerequisites}
HELIOS requires Python 3.10 or higher (recommended: Python 3.11 in Conda).
\begin{lstlisting}[language=bash]
# Clone repository
git clone https://github.com/cspark/helios.git
cd helios

# Create and activate conda environment
conda create -n helios python=3.11 -y
conda activate helios

# Install dependencies in editable mode
pip install -e ".[dev,ml,dashboard,ontology]"
\end{lstlisting}

\section{External Simulators and Control Tools}
\begin{itemize}
    \item \textbf{ANL TRACK}: Heavy-ion ray-tracing and multi-charge state tracking code executed via Wine 64-bit (\texttt{wine64 track.exe}).
    \item \textbf{LBNL IMPACT-Z}: 3D space-charge Particle-in-Cell (PIC) simulation managed through \texttt{lume-impact}.
    \item \textbf{EPICS Base}: Channel Access and pvAccess tools (\texttt{p4p}, \texttt{pyepics}) for live communication with accelerator IOCs.
\end{itemize}

\chapter{Operator Interfaces and CLI}

HELIOS provides an ergonomic command-line interface implemented via Typer and Rich.

\section{Interactive Operator Console (\texttt{helios chat})}
Launch the conversational operator assistant:
\begin{lstlisting}[language=bash]
helios chat
\end{lstlisting}
The console accepts natural language goals (e.g. \textit{``Correct horizontal orbit in MEBT and verify quench margins on SCL1 cavity 02''}), decomposing them into formal Osprey execution plans.

\section{Real-Time Web Dashboard (\texttt{helios dashboard})}
Start the FastAPI and WebSocket visualization dashboard:
\begin{lstlisting}[language=bash]
helios dashboard --host 0.0.0.0 --port 8000
\end{lstlisting}
Open \texttt{http://localhost:8000} to inspect live BPM orbits, cavity resonance states, radiation loss heatmaps, and approve human-in-the-loop plans.

\section{Commissioning Shift Reporting (\texttt{helios report})}
Compile electronic logbook reports:
\begin{lstlisting}[language=bash]
helios report --facility "RAON SCL1" --output shift_report.md
\end{lstlisting}

\section{Tamper-Evident Audit Verification (\texttt{helios audit})}
Verify SHA-256 cryptographic hash chaining across all recorded Process Variable writes:
\begin{lstlisting}[language=bash]
helios audit --verify
\end{lstlisting}

\chapter{Accelerator Physics and Beam Dynamics}

\section{Multi-Species and Multi-Charge-State (MCS) Dynamics}
Heavy-ion linacs frequently accelerate multi-charge state beams simultaneously (e.g., $^{238}\text{U}^{33+}$ and $^{238}\text{U}^{34+}$). The magnetic rigidity $B\rho$ is given by:
\begin{equation}
    B\rho = \frac{p}{q} = \frac{\gamma m_0 \beta c}{q}
\end{equation}
The module \texttt{helios.physics.species} models isotopic distributions and composite beam envelopes across charge fractions.

\section{Stripper Foil Degradation Modeling}
The module \texttt{helios.physics.stripper} simulates carbon and liquid lithium stripper foils:
\begin{itemize}
    \item Mean energy loss $\Delta E$ and stopping power $dE/dx$.
    \item Multiple Coulomb scattering angular distribution $\theta_{\text{rms}}$.
    \item Thermal degradation and foil lifetime tracking with replacement warning triggers.
\end{itemize}

\section{Beam-Based Alignment (BBA) and Dispersion-Free Steering (DFS)}
Beam offsets relative to quadrupole magnetic centers are computed via strength modulation $\Delta K_1$:
\begin{equation}
    \Delta x_{\text{BPM}} = R_{12} \cdot (\Delta K_1 L) \cdot (x_{\text{beam}} - x_{\text{quad}})
\end{equation}
Dispersion-Free Steering (DFS) solves the stacked linear system via weighted SVD:
\begin{equation}
    \min_{\Delta \theta} \left( \| R_{\text{orbit}} \Delta \theta + x_0 \|^2 + w^2 \| R_{\text{disp}} \Delta \theta + \eta_0 \|^2 \right)
\end{equation}

\section{Maximum Entropy (MENT) Phase Space Tomography}
The \texttt{helios.physics.halo} module implements iterative Maximum Entropy (MENT) reconstruction of 2D transverse phase-space $(x, x')$ distributions from 1D profile projections. The Wangler-Crump halo parameter $H$ is evaluated as:
\begin{equation}
    H = \frac{\langle r^4 \rangle}{2 \langle r^2 \rangle^2} - 1.0
\end{equation}
providing automatic warnings when beam halo risks cryomodule chamber scraping.

\chapter{Autonomous RF Conditioning and Cavity Control}

\section{Progressive RF Conditioning Sequencer}
The \texttt{AutonomousRFConditioner} automates the processing of high-power RF cavities:
\begin{itemize}
    \item Progressive forward power and pulse width ramping with vacuum interlocks ($P_{\text{vac}} < 10^{-7}\text{ mbar}$).
    \item Dynamic multipacting barrier breakthrough: delivers $+25\%$ power jumps upon detecting electron pickup spikes.
\end{itemize}

\section{Resonance Tracking and Detuning Compensation}
The \texttt{CavityResonanceTracker} provides dual-actuator control under microphonics and helium bath pressure drift:
\begin{enumerate}
    \item \textbf{Fine Fast Control}: Piezoelectric actuators compensate high-frequency phase errors ($\pm 100\text{ Hz}$).
    \item \textbf{Coarse Slow Control}: Stepper motor tuners absorb large thermal drifts and restore piezo voltage to center bias.
\end{enumerate}

\chapter{Optimization and Self-Driving Policies}

\section{Reinforcement Learning Continuous Control}
\begin{itemize}
    \item \textbf{\texttt{LinacGymEnv}}: Standard Gymnasium environment wrapping digital twin simulations and live EPICS control.
    \item \textbf{\texttt{LinacRLAgent}}: PyTorch Actor-Critic continuous policy trained for dynamic orbit correction under space charge.
    \item \textbf{\texttt{NeuralMPCController}}: Receding-horizon model predictive controller utilizing neural surrogates.
\end{itemize}

\section{Multi-Objective Pareto Optics Optimization}
The \texttt{MultiObjectiveParetoOptimizer} implements the NSGA-II non-dominated sorting genetic algorithm to simultaneously balance:
\begin{enumerate}
    \item Maximizing beam transmission ($>98\%$).
    \item Minimizing transverse envelope scalloping.
    \item Minimizing longitudinal energy spread ($\Delta p / p$).
\end{enumerate}

\chapter{Automated Fault Recovery and Self-Healing}

\section{Autonomous RF Rematching}
Upon cavity trip or phase unlock in superconducting cryomodules, \texttt{helios.healing.rf\_rematch} recalculates downstream cavity voltages and phases via normalized L-BFGS-B optimization to preserve target beam arrival time and final energy.

\section{Magnet Power Supply Recovery}
\texttt{helios.healing.magnet\_recovery} isolates tripped magnet power supplies and recalculates compensatory dipole kicks using remaining healthy steerers.

\chapter{LinkML Ontology and Knowledge Graph}

HELIOS formally defines its accelerator domain using the \textbf{LinkML} standard (\texttt{helios/ontology/helios\_schema.yaml}):
\begin{itemize}
    \item \textbf{Schema Classes}: Defines \texttt{BeamlineElement}, \texttt{MagneticElement}, \texttt{Quadrupole}, \texttt{RFCavityElement}, \texttt{BeamPositionMonitor}, and \texttt{LinacLatticeModel}.
    \item \textbf{W3C RDF / Turtle OWL Exporter}: Serializes digital twins into Linked Open Data knowledge graphs with upstream and downstream connection topology.
    \item \textbf{Topology Navigator}: The \texttt{BeamlineKnowledgeGraph} enables spatial and causal queries along the beamline (e.g., locating nearest upstream correctors).
\end{itemize}

\chapter{Digital Twin Synchronization and Security}

\section{Ensemble Kalman State Estimation}
\texttt{helios.physics.sync} implements a continuous Kalman filter assimilating live BPM telemetry to estimate unmeasured physical states, including beam energy spread and emittance.

\section{Role-Based Access Control and Watchdogs}
\begin{itemize}
    \item \textbf{RBAC Security}: HMAC-signed tokens enforce permissions across \texttt{VIEWER}, \texttt{PHYSICIST}, \texttt{SAFETY\_OFFICER}, and \texttt{ADMIN}.
    \item \textbf{MPS Heartbeat Watchdog}: Sub-second watchdog enforcing fail-safe fallback upon communication disruption.
    \item \textbf{Cryptographic Audit Trail}: SHA-256 hash chaining guaranteeing tamper evidence and regulatory compliance.
\end{itemize}

\chapter{Verification and Test Suite}

HELIOS includes an automated test suite verifying all 29 development phases:
\begin{lstlisting}[language=bash]
pytest
\end{lstlisting}
The suite validates 98 tests covering lattice physics, neural surrogates, EPICS communication, RL environments, fault recovery, LinkML schemas, and end-to-end commissioning harnesses with 100\% passing status.

\end{document}
